\begin{titlepage}
\centering
{\Large 投屏技术课教学笔记\par}
\vspace{1.0cm}
{\huge\bfseries 机器人基座模型与 VLA 经典论文路线图}\par
\vspace{0.5cm}
{\Large Zhang Xiaojun Podcast 对话陈建宇\par}
\vspace{0.35cm}
{\large 基于投屏版视频、转写稿、章节结构、视频描述与投屏关键帧整理\par}
\vspace{0.3cm}
{\large 2025-04-07\par}
\vspace{0.85cm}
\includegraphics[width=0.88\textwidth,height=0.42\textheight,keepaspectratio]{cover.jpg}\par
\vfill
\begin{tcolorbox}[width=0.92\textwidth, colback=black!2!white, colframe=black!55, sharp corners]
\textbf{视频标题}: 逐篇解析机器人基座模型和VLA经典论文 (含投屏版)——“人就是最智能的VLA”\par
\textbf{视频频道}: Zhang Xiaojun Podcast\par
\textbf{Duración del video}: 02:29:41\par
\textbf{Enlace del video}: \href{\videourl}{\nolinkurl{\videourl}}\par
\textbf{素材说明}: YouTube 未暴露可下载字幕；正文基于 faster-whisper large-v3 CPU int8 分块转写、视频描述、章节结构和投屏关键帧整理.本视频含大量投屏内容, 正文使用真实投屏帧作为视觉主线.\par
\textbf{Página del proyecto}: \href{\repourl}{\nolinkurl{\repourl}}\par
\end{tcolorbox}
\end{titlepage}

\begingroup
\small
\setlength{\parskip}{0pt}
\renewcommand{\baselinestretch}{0.92}\selectfont
\tableofcontents
\endgroup
\newpage

